\section{模型评价与推广建议}

\subsection{模型优点}

第一，模型链条完整。本文从老人数量预测出发，逐步连接服务需求、站点配置、价格补贴和情景鲁棒性评价，避免四个问题各自建模导致口径不一致。

第二，约束体系较完整。选址模型同时纳入预算、半径、容量、满意度和服务承接约束，定价模型同时纳入补贴上限、价格满意度、紧急救助免费和利润率约束。

第三，求解过程可复现。问题二采用全枚举剪枝，问题三采用临界点枚举，均适合本题 10 个小区的小规模离散优化结构，避免随机算法带来的结果波动。

第四，评价指标具有解释力。本文不只报告覆盖率，还报告容量满足率、有效服务率和服务加权满意度，能够揭示覆盖已经充分但容量仍然不足的关键矛盾。

\subsection{模型局限}

第一，同一健康状态老人被视为具有相同需求，未进一步刻画年龄、收入、家庭照护和个体偏好的差异。第二，服务站容量按总人次统一折算，没有区分助餐、护理、康复和助浴等服务对人员资质和设备的不同要求。第三，距离矩阵未反映小区内部楼栋分布、实际步行路径和老人行动能力差异。第四，满意度权重直接来自附件规则，尚未通过实际问卷或历史运营数据进行估计。

\subsection{改进方向}

后续可从三个方向改进。第一，构建分项目容量模型，将护理人员、康复设备、助浴设备和助餐配送能力分别设置约束。第二，引入随机需求或机会约束，处理季节变化、突发疾病和临时照护需求。第三，利用 GIS 路网距离和社区楼栋级人口分布数据，建立更精细的可达性和满意度模型。若进一步获得服务站运营数据，还可估计价格弹性和满意度权重，形成更接近现实运行机制的政府、服务站和老人三方协同优化模型。
